Los resultados experimentales ratifican que el comportamiento observable en arquitecturas físicas contemporáneas está fuertemente modulado por las constantes ocultas y la interacción con la jerarquía de memoria caché, más allá de la cota asintótica formal. En ordenamiento, los cuatro algoritmos verificaron la tasa $\mathcal{O}(n \log n)$, destacando la superioridad de \texttt{std::sort} (159.6 ms para $N = 10^7$ en aleatorio $D7$) debida a su baja sobrecarga interna frente al elevado costo espacial de PatienceSort ($\approx 184.6$ MB) originado por su gestión dinámica de pilas. Por su parte, en multiplicación matricial, a pesar de que Strassen ostenta una cota asintótica superior de $\mathcal{O}(n^{2.81})$ frente a $\Theta(n^3)$, el método Naive resultó más de 50 veces más rápido en $N = 256$ (11.06 ms vs. 575.37 ms), demostrando que la penalización por fallos de caché y las 18 adiciones/asignaciones de submatrices por nivel recursivo imposibilitan alcanzar el punto de cruce práctico en instancias pequeñas y medianas sin un esquema híbrido de corte.